\documentclass[12pt]{article}
\usepackage[utf8]{inputenc}
\usepackage[T1]{fontenc}
\usepackage{hyperref}
\title{Tadyawan Language Resources for K-12 Curriculum Development}
\author{Compiled by Claude Sonnet 5 (claude-sonnet-5) for Andr\'as Kornai}
\date{}
\begin{document}
\maketitle

\section{Introduction}
Tadyawan (ISO 639-3 tdy; Glottocode tady1237) is an Austronesian language of the Northern Mindoro group, spoken by Mangyan people around southern Lake Naujan in Oriental Mindoro, Philippines. Tweddell (1970) lists four dialects: Nauhan, East Aglubang, West Aglubang and Pola. Wikipedia cites about 4,200 speakers (2000). The UP Linguistics language capsule gives EGIDS 6b (threatened) and, from a 2010 census, 2,620 users in an ethnic population of 5,100. Most Tadyawans also speak Tau-buid. The UP page says little research exists on Tadyawan.

\section{Primary Educational Resources}
No primers, textbooks or school programmes were found.

\section{Linguistic and Reference Materials}
\begin{itemize}
\item Wikipedia, Tadyawan language: \url{https://en.wikipedia.org/wiki/Tadyawan_language}
\item UP Linguistics (Katig Collective) language capsule: \url{https://linguistics.upd.edu.ph/the-katig-collective/language-capsules/tadyawan}
\item Wikipedia, Northern Mindoro languages: \url{https://en.wikipedia.org/wiki/Northern_Mindoro_languages}
\end{itemize}

\section{Religious and Community Resources}
\begin{itemize}
\item Joshua Project, Tadyawan in the Philippines: \url{https://joshuaproject.net/index.php/people_groups/15169}
\item NCCA page on the Mangyan: \url{https://ncca.gov.ph/?p=7864}
\item Wikipedia, Mangyan: \url{https://en.wikipedia.org/wiki/Mangyan}
\end{itemize}

\section{Digital and Audio Resources}
None found.

\section{Curriculum Development Context}
Very little Tadyawan material exists in what was found, and nothing for schools. The sources are general reference pages. Any curriculum would need new documentation and community work, possibly drawing on materials for neighbouring Mangyan languages. I did not search for those.
\end{document}